% ======================================================================
\section{Development process and history}\label{sec:history}

The system was not designed in one step. It reached \Dp{} through fourteen phases and, in the last week, more
than sixty numbered, pre-registered rounds. This section explains the working rules, then the history in three
tables: the phases (Table~\ref{tab:eras}), the chain of shipped changes (Table~\ref{tab:shipchain}) and the main
rejected ideas (Table~\ref{tab:rejected}). Appendix~\ref{app:rounds} has one row per round.

\subsection{Working rules}\label{sec:rules}
From 13 Sept the following rules governed every change (\fp{docs/plan_robustness.md},
\fp{docs/prereg_round13_detector_push.md}):
\begin{enumerate}
  \item \textbf{Pre-register.} Before a run, write the idea, the exact rule, the data it is decided on and the pass
        bar, and commit it. Results are appended below the pre-registration, never edited into it.
  \item \textbf{Decide on development data only.} The development set (or, for audio-only rules, the held-out 415)
        decides. The test set is read after a change ships and is only reported. The author had to restate this rule
        on 1 Oct after a test-informed choice (Section~\ref{sec:disclosures}).
  \item \textbf{General rules only.} A rule must make sense for any video, run on one video at a time, and not name a
        clip or a single sound class (a rule that only ever fires on bells was rejected for this reason, Round 64).
  \item \textbf{Compare with the shipped pipeline.} Every candidate is scored as a full-pipeline variant on the same
        clips as the current shipped system, not on a proxy; a component screen may stop an idea early but cannot
        ship it.
  \item \textbf{Hits first.} Losing a few hits is acceptable only if many wrong pictures go (the fewer-pictures
        clause, Section~\ref{sec:passrules}).
  \item \textbf{Consult and record.} Major choices were reviewed by panels of independent model reviewers (Claude
        instances with assigned roles, two to five per panel, one to four rounds), whose reports are in
        \fp{docs/review/} and \fp{docs/panel_*.md}; every decision and its reason is in a dated log.
\end{enumerate}

\subsection{Phases}
\begin{longtable}{L{1.7cm}L{4.2cm}L{3.6cm}L{5.2cm}}
\caption{Project phases. Scores in each row are in that phase's own metric and are not comparable across rows.}\label{tab:eras}\\
\toprule
dates & system & scored by & outcome and why it changed \\
\midrule\endfirsthead
\toprule dates & system & scored by & outcome and why it changed \\ \midrule\endhead
\bottomrule\endfoot
06 Aug & proposal: 7 stages (audio, VLM, ASR, BEATs/PANNs, reasoning, FLUX/SDXL, evaluation) & planned: VLM describes, LLM writes reference, LLM judges & pivot from an earlier speech-illustration project \\
07 Aug & v1: PANNs (bar 0.15) $\to$ Openverse image retrieval $\to$ side panel; no gate & demo clips & needed a ``not seen'' decision \\
07 Aug--07 Sep & CLIP zero-shot visibility gate; clip-level benchmark (274 clips) & clip-level gate accuracy & CLIP gate 49.3\,\%: 71 of 103 seen clips still augmented -- ``object presence is not event visibility'' \\
08--11 Sep & PANNs + OWLv2 + per-sound VLM gate (Qwen2.5-VL-7B, ``name it, don't confirm''); SDXL & 100 clips $\times$ 3 systems, Mistral-7B judge 0--4, model-written reference & blind 2.94 $>$ caption 2.68 $>$ gated 2.41; retrieval beat SDXL by $+0.17$/$+0.18$; reference later found circular \\
12--17 Sep & v3: BEATs (replaced PANNs after 11 of 14 complaints traced to its labels), hysteresis + occlusion onsets, 7B gate with 3 votes, FLUX.1-schnell & human-grounded judge reference, dev 174 / test 100 & tie 3.17 vs 3.16; gate wins where nothing is due ($+0.26$), loses where a picture is due ($-0.24$) \\
18--21 Sep & v4: one state-of-the-art swap per stage; Qwen3.8-27B gate; FLAM and PretrainedSED detectors tried; rubric-enforced judge & judge with a cap on redundant pictures & PretrainedSED rejected three times; 13 silent bugs found (amendment 3), all earlier judge rows kept as history only \\
22--23 Sep & per-sound gold (139 clips); BEATs $\cup$ FlexSED (bar 0.8), FlexSED veto 0.3, PANNs veto 0.05 & per-sound F1, then viewer cost (amendment 9) & F1 difference null ($+0.028$); wrong pictures and precision significant \\
24--27 Sep & monotone onsets, no length cap; final September test table (amendment 21); Qwen-Image frozen as picture setup & Holm-corrected table on TEST-60 & ``Holm wins, F1 null'' (Section~\ref{sec:sept-table}) \\
27 Sep & detector rounds 1--3: cascade, PE-A-Frame, Qwen3-Omni as verifier & DEV-49 + AudioSet 280 / 415 & none passed; detection declared closed, reopened next day \\
28 Sep & rounds 4--12 on AudioSet sets: EAT, Dasheng, DASM, SAM-Audio, 17 ideas, local contrast, DASM rescue, masker tests & cost on AudioSet labels & all failed on the 415; round 12 showed the AudioSet cost cannot track the task: harness retired \\
28--30 Sep & rounds 13--14: listener rescue (Qwen3-Omni, Audio Flamingo Next) + DASM vote; tagger clips added & DEV-49, then merged DEV & first gains in hits (TO1+F7F8) \\
30 Sep--01 Oct & ship chain SHIP1--SHIP8 on merged DEV/TEST & merged DEV, pass rules & DEV 18/51/3.690 $\to$ 28/21/2.282 (Table~\ref{tab:shipchain}) \\
01--02 Oct & rounds 46--66: display rules, witness rule, scene check, many gate ideas & merged DEV, held-out 415 & \Dp{}: DEV 29/15/2.056; frozen 2 Oct 01:20 UTC \\
\end{longtable}

\subsection{The shipped chain}\label{sec:shipchain}
Table~\ref{tab:shipchain} is the backbone of the history: every change that entered \Dp{}, in order, with its
development and test numbers on the merged sets. Development numbers use the corrected gold (58 needed sounds);
test numbers are the reads made after each change shipped. Sources: \fp{docs/review/improvements_log_2026-09-30.md},
\fp{docs/review/detector_final_2026-10-02.md}, \fp{docs/prereg_round13_detector_push.md}.

{\small
\begin{longtable}{L{2.6cm}L{6.1cm}L{2.5cm}L{3.0cm}}
\caption{Changes that shipped, from the old pipeline to \Dp. Hits / wrong / cost; DEV hits out of 58, TEST out of 65.
The $p$ in the TEST column is one-sided, against the old pipeline B0r for SHIP5--SHIP8 and against B for D and \Dp.}\label{tab:shipchain}\\
\toprule
step (name in logs) & what it adds & merged DEV & merged TEST \\
\midrule\endfirsthead
\toprule step & what it adds & merged DEV & merged TEST \\ \midrule\endhead
\bottomrule\endfoot
old pipeline (B0r) & September test-table system: BEATs $\cup$ FlexSED, FlexSED and PANNs vetoes, Qwen3.8 gate & 18 / 51 / 3.690 & 21 / 40 / 2.909 \\
TO1+F7F8 (R14) & listener rescue of weak FlexSED runs (tiered: Qwen alone above 0.6, both listeners below), one rescue per family, mirror veto with listener keep, DASM vote on rescues & 26 / 45 / 3.070 & 22 / 38 / 2.818 \\
+ N2b, DR2, K-V4 (R16--21) & masked-weak veto; DASM rescue of new families; a listener keep must also be named in an inventory & 27 / 38 / 2.817 & 23 / 36 / 2.727 \\
+ DV (R28) & drop a sound DASM never hears in the clip & 27 / 31 / 2.620 & 23 / 34 / 2.682 \\
+ BTP (R30) & start a late picture at an earlier weak run of the same sound & 28 / 30 / 2.535 & 23 / 33 / 2.659 ($p$ 0.033) \\
+ CONT (R31) & drop a picture that is only a later piece of a sound already going on & 28 / 25 / 2.394 & 24 / 32 / 2.591 ($p$ 0.025) \\
+ FineLAP veto (R31) & a 2025 frame-level model double-checks rescued sounds & 28 / 24 / 2.366 & 24 / 31 / 2.568 ($p$ 0.015) \\
+ K4A-D = SHIP8 (R35) & drop a picture neither listener names, unless DASM clearly hears it & 28 / 21 / 2.282 & 23 / 29 / 2.568 ($p$ 0.017) \\
+ merge gap 2.5\,s & continuing sounds with short breaks no longer drew a second picture & 28 / 20 / 2.254 & 23 / 28 / 2.545 \\
+ grouping $\leq$8\,s (R47) & Qwen3-Omni: same continuing sound or a new event? & 28 / 18 / 2.197 & 22 / 26 / 2.545 \\
+ shortest sound 0.3\,s (R48) & short real sounds (a laugh burst) were thrown away & 29 / 18 / 2.141 & 22 / 26 / 2.545 \\
+ depict-event (R57) & drop a picture whose event is visibly happening & 29 / 17 $\to$ 18$^{a}$ / 2.141 & 22 / 26 / 2.545 \\
= \textbf{B} (SHIP8+MD3) & base of the last rounds & \textbf{29 / 18 / 2.141} & \textbf{22 / 26 / 2.545} \\
+ witness rule + scene check (R53+R60) = D & weak sounds need DASM, both listeners, or one listener and a plausible scene & 29 / 14 / 2.028$^{b}$ & 24 / 23 / 2.386 \\
+ logit readout (R60L) = \textbf{\Dp} & scene check read as a probability, no cut-off answers & \textbf{29 / 15 / 2.056} & \textbf{24 / 24 / 2.409} ($p$ 0.034 vs B) \\
\end{longtable}
}
\noindent\footnotesize $^{a}$ Round 57 scored 29/17 with 26 listener answers missing from a cache; after the cache
was filled (1 Oct evening) the same system scores 29/18/2.141, so its net development effect is zero.
$^{b}$ D's 2.028 contains one gain caused by a cut-off answer (``based on the'') that removed a wrong \emph{Glass}
picture; the logit readout of \Dp{} cannot be cut off, so that picture returns (+1 wrong). The pass bar of Round 60L
was set against B, not against D, for this reason (prereg lines 4504--4542).
\normalsize

\paragraph{What the table shows.} From the old pipeline to \Dp{} the development set gained 11 hits and lost 36 wrong
pictures (cost $3.690 \to 2.056$). The test set gained 3 hits and lost 16 wrong pictures ($2.909 \to 2.409$). About a
third of the development gain in cost carried over to the test set; with well over a hundred development variants
tried, some of the development gain is expected to be selection luck (``winner's curse''), and the test numbers are
the better estimate of what a new video would see.

\subsection{What was tried and rejected}\label{sec:rejected}
Most ideas failed. Table~\ref{tab:rejected} groups the main ones by the question they tried to answer, with the reason
from their own pre-registered test. All were decided on the development set or the held-out 415.

{\small
\begin{longtable}{L{3.4cm}L{5.3cm}L{5.8cm}}
\caption{Main rejected ideas (selection). Hits / wrong / cost on merged DEV unless stated.}\label{tab:rejected}\\
\toprule
idea (round) & what it did & why it failed \\
\midrule\endfirsthead
\toprule idea & what it did & why it failed \\ \midrule\endhead
\bottomrule\endfoot
\multicolumn{3}{l}{\emph{Hear more sounds}} \\
other detectors (Sep, rounds 1--8) & FLAM, PretrainedSED (5 backbones), AST, CED, SSLAM, EAT, Dasheng, PE-A-Frame, CLAP verifier, tagger ensembles & recall rose only with many more false spans; PretrainedSED lost three times; ensembles mis-scaled \\
source separation (Sep, R7) & Demucs, SAM-Audio in front of BEATs & separated audio is out of the detectors' training distribution: recall under speech fell 9.5\,\% $\to$ 4.8\,\%; SAM-Audio failed its quality check (music never removed) \\
listener as verifier (R13) & ask Qwen3-Omni yes/no, multiple choice, open list on every weak candidate & finds 5--8 dropped real sounds but adds about 3 wrong pictures per hit \\
newer listeners (R31--34) & MOSS-Audio, SpotSound, Kimi, Step-Audio & 3--10 wrong per extra hit; SpotSound said ``Yes'' to everything \\
EXPECT chain (R40--42b) & the scene VLM or a listener names sounds expected off screen; FineLAP times them; DASM confirms & best DEV 31/26/2.254 (+3 hits), but TEST 25/35/2.614 (worse); the gate cannot refuse a wrong off-screen name \\
short twins, band lists (R49, R56) & add weak or short detections both models half-hear & BANDLIST 0 hits, +6 wrong; TWIN-SHORT 27/36/2.761 \\
tagger ensemble (R63) & calibrated EAT/SSLAM ensemble on BEATs windows & on the 415: 33\,\% fewer spans but recall 0.152 vs 0.197 \\
\multicolumn{3}{l}{\emph{Remove wrong pictures}} \\
stricter witnesses (R53 variants) & require DASM or both listeners for every weak sound & 27/13/2.113: $-5$ wrong but 2 real hits lost (fails the 3:1 clause) \\
named veto, relabel gate (R39--41) & drop or rename a picture when the gate or the listeners name a different sound & real hits lost for no cost gain \\
held-out family prior (R41) & drop families a detector gets wrong $>70\,\%$ of the time on the 415 & every affected picture had been heard by a listener: no change \\
\multicolumn{3}{l}{\emph{Better on-screen decision}} \\
other gate signals (Sep, am.\ 14--20) & OWLv2 or SAM 3 vote, video input, synchrony, ``any stretch'' rule & each removed exactly as many needed sounds as wrong pictures it saved (break-even at $\beta = 2$) \\
SSL-SaN, SigLIP, BOX (R30--38) & audio-visual localiser; picture-vs-frame similarity; VLM draws a box on the source and checks the crop & each un-silenced real visible sounds; BOX-2 in the full pipeline 29/26/2.366 \\
synchrony (R38) & Synchformer: does on-screen motion move with the sound? & as a veto, +3 needed kept but 5 visible calls lost; as an ``add seen'' vote, on the bar but changed no shipped picture \\
human-style chain (HUMAN, R50, R66) & name the source, find it, zoom, ask ``is it making the sound now?'' & the VLM names ``nothing'' on most stretches; NAME-ALL passed its gate screen but lost 3 real hits in the pipeline (26/12/2.141) \\
sign/event and logit gate (R54, R62) & ask for visible effects; read the gate as a logit margin & letter-position bias or ``no'' habit; best AUROC 0.647, below the 0.65 bar \\
look-alike reasoning (R61) & ``no train here, so it is the washing machine'' & works on the laundromat, but also explains away real sounds (a crying sound ``explained'' by a visible parrot) \\
bell rule (R64) & concealed actions (a bell in a tower) count as not visible & passed DEV (30/14/1.972) but fires only on bells: not general, not shipped \\
\multicolumn{3}{l}{\emph{Timing and grouping}} \\
change-point onsets (R33) & SEBB onset post-processing & 7 hits lost \\
earliest onset (R38), repeat rules (R31, R35) & start at the first detector to hear a sound; require silence before a repeat & hits lost or no net gain \\
pause questions (R51, R55) & ask whether a 2-s gap is a pause or a new sound & neither the listener nor the detectors can tell a pause from a dropout \\
context $\pm$5\,s (R59) & give the listener more context & it names what it \emph{sees}, not what it hears \\
\end{longtable}
}

\paragraph{The main lesson.} Two patterns recur in this table. First, the missing sounds can often be found -- a
listener asked about the whole clip names the hammer, the train, the bell -- but every way of adding them brings two
to four wrong pictures per hit, because nothing downstream can tell a right off-screen name from a wrong one. Second,
every change to the gate moves along the same trade-off line: it silences more on-screen sounds only by silencing
about as many needed ones. Section~\ref{sec:analysis} measures both.

\subsection{Pictures}\label{sec:pic-history}
The picture model changed four times: Openverse retrieval (v1), SDXL and SDXL-Turbo, FLUX.1-schnell (chosen over
PixArt-$\Sigma$ on twelve real prompts, 13 Sept), and Qwen-Image-2512 (frozen 25 Sept, shipped 28 Sept). On 9 Sept a
retrieval-versus-SDXL ablation on the same 100 clips favoured retrieval by $+0.17$ (gated) and $+0.18$ (blind) judge
points, but 55\,\% of retrieved images carried a non-commercial or no-derivatives licence and 595 uses drew on only 60
distinct images, so generation was kept. The decisive picture test was a blind \emph{glance test} by the author on 25
Sept: 54 sounds from 50 never-annotated clips, each picture shown for 1.5\,s at 384\,px with no name, and the author
typed what the sound was. FLUX was recognised 14/54 times; Qwen-Image with the same text 26/54
($+0.222\ [+0.074, +0.370]$); Qwen-Image with a richer subject text 32/54, but that text invented an object (a thud
drawn as a door) and was rejected. One of the 26 Qwen-Image pictures was later found to send a false message (a
telephone bell drawn as a desk bell), so the $+22$-point result is reported as ``not clean''. A planned sealed
confirmation sitting was cancelled on 28 Sept and replaced by the check-and-redraw loop of Section~\ref{sec:pictures}.
A wording test the same day kept the hand-written subject table: 7 of 7 known cases right, against 3 of 7 for two
generic VLM wording methods (\fp{docs/picture_sense_test_2026-09-28.md}).

\subsection{Evaluation history: the judge era}\label{sec:judge-era}
The proposal planned an automatic evaluation: a VLM describes the augmented video, an LLM writes a reference of what
the viewer should learn, and an independent LLM scores the description against the reference on a 0--4 scale. It
was built and run, and it failed in instructive ways:
\begin{enumerate}
  \item \textbf{A reference written by the system's own models is circular.} The first reference was written by the
        system's VLM from the system's detections; it rewarded whichever system repeated the detector (blind 2.94 $>$
        caption 2.68 $>$ gated 2.41 on 9 Sept). A second judge of another family agreed with the first
        (quadratic $\kappa = 0.753$), which shows only that two judges can share a bias.
  \item \textbf{A reference written by independent models cannot decide visibility.} A reference built by an audio
        LLM, CLAP, a different VLM and a different LLM agreed with the human ``nothing is missing'' label on 55 of 100
        clips -- chance. An audio-visual model asked with and without the soundtrack reached 49.5\,\% balanced
        accuracy (bar 67.7\,\%); a ``gap-closing'' score had AUROC 0.529 (bar 0.75).
  \item \textbf{The rubric was asymmetric.} A withheld picture scored 0 where one was due, a redundant one about 3
        where none was due, so the rubric punished staying silent -- the very behaviour the gate exists for. A capped
        rubric (19 Sept) and a human-grounded reference fixed the direction but not the noise.
  \item \textbf{Judges cannot tell a picture from a text tag.} The best judge (Gemma-4-31B seeing the pictures
        directly, passing all three trust checks against the author's labels) scored pictures and text tags alike
        ($-0.04\ [-0.22, +0.14]$), because it reads a tag as the label. Six automatic picture checkers also failed
        calibration against the author (for example $\kappa = 0.21$ to $0.45$ against bars of 0.5).
\end{enumerate}
A five-reviewer scoring panel on 19 Sept moved the primary score to per-sound matching against human labels
(Section~\ref{sec:task}) and kept the judge as a secondary instrument. The answer to RQ4 is built on this history.
